\subsubsection{The long-orbit row bound}\label{sec:inert-long-orbit}
\hypertarget{source-inert-long-orbit}{}

Among the inert orbits, a long Frobenius orbit allows the marked jumps to be cut into shorter arcs. The colength can be bounded using those arcs only if the correction at each cut is controlled integrally.

\begin{established}[Case B and the cyclic signed jump windows]\label{hyp:long-orbit-data}
Let \(\ell<A_{\rm sec}\) be an ordinary inert prime whose Frobenius
has order \(o_\ell>4L_w\) in the real cyclotomic quotient
(case B of Section~\ref{sec:cut-length-rows}). Write \(F=E_\lambda\) for the
unramified completion of the auxiliary field and
\[
 d=[F:\Q_\ell]=n_\lambda.
\]

Work over the integers \(O\) of a local coefficient field splitting \(F\).
Write \(T_2\) for the marked dual Tate lattice, of rank \(r\), and
\(T_1\) for a lattice in a subrepresentation of a finite sum of the
twisted duals from the same cut. The maps to be estimated, including the marked matrix row, are precisely the indicated integral Galois surjections
\[
 f:T_1\twoheadrightarrow T_2,\qquad \operatorname{rank}_O T_2=r .
\]

The source and target are crystalline with central jumps \(|a|\le\eta_0\) and one outer signed interval of width \(S_0\) at each Frobenius embedding. The small letters have height at most \(G_P\) and all letters have height at most \(H\), in the notation of Section~\ref{sec:cut-length-rows}.
\end{established}
\begin{proof}
The ordinary inert branch excludes all primes ramifying in the field or in the rows. Conjugating its marked row preserves the coefficient lattice and takes its dual Tate jumps to the reflected signed intervals. Apply Lemma~\ref{lem:marked-row-column-lattices} to obtain the displayed surjection. Lemma~\ref{lem:marked-signed-jumps} gives the central and outer bands before passage to its quotient. Both rational representations are crystalline with the
local jump windows from the cut. Enumerate the embeddings of \(F\)
cyclically by Frobenius. The residue characteristic is \(\ell\ge11\)
in this inert branch.

Use \(b_0,w_0,S_0,\eta_0\) for the local size parameters, with
\(S_0,\eta_0\ge1\), and \(s_i\) for the size at position \(i\).
The source jumps at \(i\) lie in \([-\eta_0,\eta_0]\) together
with either \([s_i,s_i+S_0]\) or its negative. The larger range is
\[
 b_0\le s_i\le b_0+\eta_0.
\]

At every other position,
\(s_i\le w_0\). Put \(g=w_0+S_0+\eta_0\) and set
\[
 h=\begin{cases}b_0+\eta_0+S_0,&\text{if some position has larger jumps},\\
                  g,&\text{otherwise}.\end{cases}
\]
Thus the positions with smaller jumps have jumps in \([-g,g]\) and all positions
have jumps in \([-h,h]\). The target satisfies the same bounds by
strict exactness of the rational crystalline filtration.
The two upper endpoints in Section~\ref{sec:cut-length-rows} bound \(b_0+\eta_0+S_0\le H\)
and \(g\le G_P\), with \(H,G_P\) as in Section~\ref{sec:cut-length-rows}.
The length of the cycle here is the unramified local degree, including
the quadratic components of Frobenius.
\end{proof}

\begin{auxlemma}[Spacing of the larger jumps]\label{lem:long-orbit-spacing}
On the unramified cycle of Lemma~\ref{hyp:long-orbit-data}, let
\(J\) be the positions with larger jumps and \(L_w\) the length
parameter for \(T_2\) in Lemma~\ref{lem:second-set-progressions}.
Write \(\theta_\ell\in\mathbb Z/D\mathbb Z\) for the action of
arithmetic Frobenius on \(C_e\). If
\(o_\ell=\operatorname{ord}(\theta_\ell)>4L_w\), then
\[
 \#(J\cap I)\le2\qquad\text{for every interval of }L_w
                  \text{ consecutive positions }I .
\]
If \(1\le j_\ell\) and \(2j_\ell+2\le L_w\), the cycle admits
cuts with \(j_\ell\) smaller positions following each cut, at most
two larger positions per arc, and maximum arc length \(3j_\ell+1\).
\end{auxlemma}
\begin{proof}
In \(\operatorname{Gal}(C_e/\mathbb Q)\), the step \(\theta_\ell\)
is the class of \(\ell\bmod 2^{e+2}\) up to sign, expressed in the
generator 5. Thus
\[
   o_\ell=\min\{a\ge1:\ell^a\equiv\pm1\pmod{2^{e+2}}\},
    \qquad o_\ell\mid [E_\lambda:\mathbb Q_\ell] .
\]
The divisibility holds because the order of the full Frobenius
equals the unramified local degree and restriction preserves powers.

By Lemma~\ref{hyp:translate-weights} the larger places are the
translated labels \(e_1(s+u-t)\) and their complex conjugates
at the two embeddings of a real place when the coordinates are
centred at the selected top \(t\). Restriction to the real
field \(C_e\) kills complex conjugation. Let \(w\) be the
coordinate of the translate, equal to \(\rho(u-t)\) with this
origin or \(\rho(u)\) before recentering. With \(s_i\) the integer representative of the larger label at
position \(i\), \eqref{eq:dyadic-coordinate-orders} gives
\begin{equation}\label{eq:translated-dyadic-visits}
 x_0+i\theta_\ell=w+t_0 s_i\pmod D,\qquad
                  0\le 4L_w s_i<n_0\quad(i\in J).
\end{equation}

Here \(x_0\) is the starting coordinate. Reversing the orientation
replaces \(\theta_\ell\) by its negative and preserves its order.
If \(a<b<c\) in one interval of length \(L_w\) are larger
positions, their coordinates are distinct by the order bound.
Put \(u'=b-a\), \(v'=c-a\). Eliminating the step gives
\[
 t_0\big(v'(s_b-s_a)-u'(s_c-s_a)\big)\equiv0\pmod D,
 \qquad
 v'(s_b-s_a)-u'(s_c-s_a)\equiv0\pmod{n_0}.
\]

The second implication uses \(D=t_0n_0\). The four terms have
absolute value less than \(n_0/4\) by \eqref{eq:translated-dyadic-visits},
so the last congruence is an integer equality between three
distinct selected labels. The sphere construction in
Lemma~\ref{lem:second-set-progressions} excludes this equality.
Taking integer lifts of the positions includes intervals crossing
the end of the cycle.

Let \(j=j_\ell\). The numerical host choice allows some exponent
\(j'\) with \(A_{\rm sec}\le(\ell-1)^{j'}\) and \(2j'+2\le L_w\).
Here \(\ell<A_{\rm sec}\), so the least exponent \(j\) is at least two
and satisfies
\[
 2j+2\le L_w.
\]
Thus the local cycle has \(d>2j\).
Join successive positions with larger jumps when their forward gap is at most
\(j\). A chain of three would occupy at most \(2j+1\) consecutive
positions and is excluded. The clusters thus have one or two
positions, with at least \(j\) intervening positions with smaller jumps between
clusters. Cut just after the last position with larger jumps of each cluster,
using that position as the cut index \(c\). Each of \(c+1,\dots,c+j\)
then has smaller jumps.

If the next cluster starts at forward distance \(a\),
insert cuts at the multiples of \(j\) up to \(a-j-1\), measured
from the current cut. Include the original cut at distance zero. If
\(k=\lfloor(a-j-1)/j\rfloor\), the remaining gap is
\[
 a-kj=j+1+\bigl((a-j-1)\bmod j\bigr)\in[j+1,2j].
\] Adding the span of the next
cluster gives an arc length at most \(3j\). Each arc is taken
half-open, with its terminal cut included.

Each such arc contains at most two
positions with larger jumps. When there are no positions with larger jumps, partition the
cycle into arcs of length \(j\) with the last remainder added to one
arc. Each arc has length at least \(j\), since inserted cuts are
spaced by \(j\) and each gap to a cluster of positions with larger jumps starts with its
\(j\)-letter buffer of positions with smaller jumps. The maximum length \(A\) satisfies
\(j\le A\le3j+1\).
\end{proof}

\paragraph{The exponent in the Wach lifting estimate.}\label{par:inert-cut-exponents}
The correction at a cut in Paragraph~\ref{par:wach-cut-correction} and the integral level tests of Proposition~\ref{prop:wach-arc-membership} apply to the
cyclotomic Wach modules of a surjection of crystalline lattices over
an unramified field. Its arc hypotheses are a cyclic partition with
\(j\) positions with smaller jumps after each cut, at most two positions with larger jumps per
arc, and specified size inequalities when their signs differ.
Lemma~\ref{lem:long-orbit-spacing} supplies the partition.

We now compute the exponents for the symmetric jump boxes of
Lemma~\ref{hyp:long-orbit-data}.
Put \(q_0=\ell-1\), \(D_0=q_0^{j+1}\), and
\[
 \begin{split}
 M&=\left\lfloor\frac{h+2g}{q_0}\right\rfloor+1,\\
 \sigma&=\left\lceil 1+\frac{\ell}{q_0^2}
                      \left(3g+\frac{M}{\ell^{j-1}}\right)\right\rceil,\\
 H_i&=\begin{cases}2h,&i\text{ with larger jumps},\\2g,&i\text{ with smaller jumps},\end{cases}\\
 c_H&=\left\lceil 1+\log_\ell(1+\max_i H_i)\right\rceil,\\
 S(c)&=\sum_{a\ge0}\frac{H_{c+a+1}+c_H+3}{q_0^a},\\
 R&=\max_c\left(c_H+\left\lfloor\frac{\ell}{q_0^2}
           \min\{\lfloor S(c)\rfloor,\max_i H_i\}\right\rfloor\right).
 \end{split}
\]
The initial exponent is \(a_0=2\sigma+R+2gA\). Define \(\kappa\)
as the maximum of
\[
       A+\sum_{u\in J\setminus\{0\}}(A+1+v_\ell(u))
\]
over unions \(J\) of three intervals of at most \(3g\) consecutive
integers in \([-2h,2h]\).

In Theorem E the inputs for the source and
target can both be taken to be \(H_i\). If the positions with larger jumps of
opposite sign have magnitudes \(x\) at backward offset \(s_1\) and
\(y\) at offset \(s_0<s_1\), the required comparisons are
\(y<\ell x\) when the earlier position in the cycle is negative,
and \(2y\le\ell x\) and \(2a_0<x\) when it is positive. With these
inputs the theorem lifts every integral finite class of \(T_2\) after
multiplication by \(\ell^e\), where
\begin{equation}\label{eq:long-orbit-exponent}
                 e=2R+4\sigma+2a_0+2\kappa.
\end{equation}

For the proof with cyclotomic evaluation coordinates, extend the
coefficient field to contain the primitive \(\ell^d\)-th roots of
unity. This extends the coefficients while keeping \(F\) and its
Frobenius cycle fixed. The embeddings over the larger coefficient
field are the same after composition, since the original field
already splits \(F\). The finite cohomology group of a lattice
is the preimage of the rational finite group.

A basis of the
extended coefficient ring over \(O\) identifies its continuous
cochains with finitely many copies of the original complex.
Rationalization and the crystalline period comparison also commute
with this finite direct sum. Thus integral finite cohomology, its
map and the cokernel extend by tensor product. For the cokernel \(C\) and extended coefficient DVR \(O'\), faithful flatness gives
\[
 \ell^e(C\otimes_O O')=0\quad\Longrightarrow\quad\ell^e C=0.
\]


The geometric estimate for \(R\) exhibits the dependence on the
buffer of positions with smaller jumps. The first \(j\) terms of the sum at each cut have
\(H_{c+a+1}=2g\), hence
\[
 S(c)\le \frac{q_0}{q_0-1}
          \left(2g+c_H+3+\frac{2h}{q_0^j}\right).
\]
Since \(\ell\ge11\) and \(h\ge1\), we have
\(c_H\le3+\log_\ell h\) and
\[
  R\le8g+3h/D_0+2\log_\ell h,
  \qquad \sigma\le3g+2h/D_0.
\]
Indeed \(\ell/(q_0(q_0-1))\le11/90\) and
\(2\ell/(q_0-1)\le22/9\) give the first inequality.

For the second, \(\ell/q_0^2\le11/100\), and the coefficient
of \(h+2g\) is at most \((11/10)/D_0\) for \(j\ge1\).
Choose the common constant \(C_2\ge3200\) in the size condition
\(b_0\ge C_2(j+1)(g+1)\). This condition holds uniformly for
the rounded local word. In fact \(b_0\ge b-(t+1)/2\ge b/2\)
at large scale, \(g+1\le G_P+1\le11b'\), and
\(j+1\le3+\log A_{\rm sec}=O(\log x)\). Hence
\[
 C_2(j+1)(g+1)/b_0=O(C_2\log x\, b'/b),
\]
which tends
to zero since
\[
 C_2\log x\,b'/b=C_2x\log x\exp(-x^{\varepsilon/3})\longrightarrow0.
\]

If there are positions with larger jumps,
\(g\le b_0\), \(h\le2b_0\), and both magnitudes are between
\(b_0\) and \(b_0+\eta_0\). Therefore \(2y\le4b_0<\ell x\).
The bounds on \(\sigma,R\) in the expression \(a_0=2\sigma+R+2gA\) give
\[
 2a_0\le40g(j+1)+14h/D_0+4\log_\ell h<b_0.
\]
For example, the three terms are bounded by \(0.04b_0\),
\(0.28b_0\), and \(4\sqrt{2b_0}\) using \(b_0\ge2000\) and
\(\log_\ell h\le\sqrt h\). The two mixed-sign comparisons follow.
When all positions have smaller jumps these comparisons are vacuous.

To bound \(\kappa\), each interval with at most \(3g\) integers
contains at most \(3g/\ell^v+1\) nonzero multiples of \(\ell^v\),
and such multiples require \(v\le\log_\ell(2h)\). Summing over
the three intervals gives
\(\kappa\le25gA+3\log_\ell h\) since \(A,g\ge1\).
Using \eqref{eq:long-orbit-exponent},
\[
 e\le110gA+28h/D_0+14\log_\ell h
   \le C_2\left(h/D_0+(g+1)(j+1)+\log_\ell h\right).
\]
In particular the high jump bound carries the factor
\((\ell-1)^{-j-1}\) even when the cycle has many clusters. The
exponents use a maximum over the arcs. The factor for the host
degree enters only in the dimension of finite cohomology.

\paragraph{Correcting the lattice at a Frobenius cut.}\label{par:wach-cut-correction}
In the ordinary cyclotomic variable \(\pi\), let \(\Delta\)
be the tame group of order \(\ell-1\) and take
\(X=-\prod_{\delta\in\Delta}\delta(\pi)\).
\begin{lemma}[The tame-invariant parameter]
The tame-invariant coefficient ring is \(O[[X]]\). Put \(q=X+\ell\)
and \(p_s=\phi^s(q)\). Up to a unit, \(p_s\) is a distinguished
polynomial of degree \(\ell^s\) with constant term associated to
\(\ell\), and for each \(j\ge1\),
\[
 (\phi^j X)=(X)\prod_{0\le s<j}(p_s)^{\ell-1}.
\]
\end{lemma}
\begin{proof}
The defining product is a unit times \(\pi^{\ell-1}\). On the invariant ring
write \(\mathcal R=O[[X]]\). Indeed the action on degree \(k\) in the \(\pi\)-adic
filtration is by the \(k\)-th power of the tame character.
An invariant series can only have its first nonzero term in
degree a multiple of \(\ell-1\). Successively subtracting
a constant multiple of the corresponding power of \(X\)
expresses it as a series in \(X\).
At a primitive \(\ell\)-th root of unity minus one,
\(X\) evaluates to \(-\ell\), by the product formula for
all primitive \(\ell\)-th roots. Thus \(q=X+\ell\) is
associated to the usual cyclotomic factor in \(O[[\pi]]\)
by Weierstrass preparation.

As a result
\[
 \phi(X)=uXq^{\ell-1},\qquad
 \gamma(X)=\omega_X X\pmod{X^2},\qquad
 v_\ell(\omega_X-1)=1 ,
\]
where \(u\) is a tame-invariant unit and
\(\omega_X=\chi(\gamma)^{\ell-1}\).
Here \(\phi\) acts on coefficients in
\(\mathcal R\) and \(\Phi\) is the corresponding semilinear operator
between letters. Put \(p_s=\phi^s(q)\). Iteration gives
\(\phi^j(X)\) associated to \(X\prod_{s<j}p_s^{\ell-1}\), and
\(p_s\) lies in \((\ell,X^{\ell^s})\). The latter follows by
reducing the Frobenius formula modulo \(\ell\). By Weierstrass
preparation, \(p_s\) is a unit times a monic polynomial of degree
\(\ell^s\) whose other coefficients are divisible by \(\ell\)
and whose constant coefficient is associated to \(\ell\).
Different \(s\) have different cyclotomic zero sets and their
distinguished polynomials are relatively prime over
\(\operatorname{Frac}(O)\).

\end{proof}
For the free Wach lattices \(N_i\), use the height bounds
\(q^h\Phi N_i\subset N_{i+1}\) everywhere and
\(q^g\Phi N_i\subset N_{i+1}\) in the first \(j\) steps after a
cut. The cyclotomic action is trivial on \(N_i/XN_i\). Write
\(P_c=\mathcal R\Phi^j N_c,\ D=X^M\prod_{s<j}p_s^{3g}\).

\begin{lemma}[Correction at a cut]
Use the chosen positive integers \(M,j\) and nonnegative integer
\(\sigma\) with \((\ell-1)M>h,\ (\ell-1)M\ge3g\) and
\((\ell-1)^2\sigma/\ell-3g\ge M/\ell^{j-1}\). In the odd-prime
range of the cut, each \(z\in XN_c\) has a correction \(\tau\in P_c\)
satisfying
\[
 \ell^\sigma\Phi^j z-(\gamma-1)\tau\in DP_c\subset X^M N_{c+j}.
\]
Its class modulo \(DP_c\) is uniquely determined by the two
congruences for \(v,\tau\) below, is additive in \(z\), and is
compatible with maps of the Wach lattices.
\end{lemma}
\begin{proof}
The height bounds and the parameter formula give
\[
 \Phi(X^M N_i)\subset q^{(\ell-1)M-h}X^MN_{i+1},\qquad
 \Phi^j(X^MN_c)\subset DP_c\subset X^MN_{c+j}.
\]
The first inclusion yields contraction of \(\Phi\) on the cyclic
product of these deep lattices. For the second use
\[
 (\prod_{s<j}p_s^g)P_c\subset N_{c+j},
\]
by iteration.
Frobenius sends a basis to an independent family after inverting
the cyclotomic factors by the two-sided height comparison. It thus
gives a basis of \(P_c\). The image of \(D\) modulo a uniformizer
is a nonzero power series. Componentwise reduction in this basis
then shows that \(P_c/DP_c\) has no \(\ell\)-torsion.
Both lattices are \(\gamma\)-stable by commutation with \(\Phi\)
and stability of the ideals \((X),(p_s)\).
Given \(z\in XN_c\), first solve the equation for
\(\gamma-1\) modulo \(X^M\).

On degree \(r\) it acts by
\(\omega_X^r-1\). With
\(A_r=\sum_{k=1}^r(1+v_\ell(k))\) and \(K=A_{M-1}\), there are
\(d_r\in N_c\) such that
\[
 v=\sum_{1\le r<M}\ell^{K-A_r}X^r d_r,\qquad
 (\gamma-1)v-\ell^K z\in X^M N_c.
\]
To construct \(v\), at degree \(r\) multiply the previous
approximation by \(\ell^{1+v_\ell(r)}\) and subtract an
\(X^r\)-term given by the inverse of the unit part of
\(\omega_X^r-1\) times the degree-\(r\) residual coefficient.
This eliminates the residual at that degree. Legendre's formula
gives \(A_r\le\ell r/(\ell-1)\).

For \(1\le r<M\) one can divide
\(\ell^\sigma(\phi^j X)^r\) by \(\ell^{A_r}\) modulo \(D\).
For \(A_r\le\sigma\) this is literal division. Otherwise let
\(d_j=\ell^{j-1}\) and
\(m=(\ell-1)r-3g-(A_r-\sigma)\). From the bound on \(A_r\)
and the definition of \(\sigma\),
\[
 m\ge\left(\frac{(\ell-1)^2}{\ell}-1\right)A_r-3g+\sigma
   \ge\frac{(\ell-1)^2}{\ell}\sigma-3g\ge M/d_j.
\]

The binomial expansion of \(p_{j-1}^{A_r-\sigma+m}\), using
\(p_{j-1}\in(\ell,X^{d_j})\), lies in
\((\ell^{A_r-\sigma},X^{d_j m})\).
After multiplying by the other factors in \((\phi^j X)^r\),
the second summand is divisible by \(D\). Choose \(b_r\in
\mathcal R\) with
\(\ell^{A_r} b_r\equiv\ell^\sigma(\phi^jX)^r\pmod{D}\).
Then
\[
 \tau=\sum_{1\le r<M} b_r\Phi^j d_r\in P_c,\qquad
 \ell^K\tau-\ell^\sigma\Phi^j v\in DP_c .
\]
Commuting \(\gamma\) past \(\Phi^j\) and using the equation for
\(v\) shows that
\[
 \ell^K(\ell^\sigma\Phi^j z-(\gamma-1)\tau)\in DP_c.
\]

Saturation removes \(\ell^K\), and
\[
 DP_c\subset X^M N_{c+j}
\]
gives the required correction. If two corrections solve the
two displayed congruences for the same \(z\), take their
difference and bring their powers of \(\ell\) to a common power.
The associated difference of the \(v\)'s is in \(XN_c\) and
has \(\gamma-1\) image in \(X^MN_c\).

If \(v'\) denotes that difference, reduction in each successive degree gives
\[
 (\omega_X^r-1)[v']_r=0\ \Longrightarrow\ [v']_r=0
 \qquad(1\le r<M).
\]
The scalar is nonzero on a free graded piece, so induction gives \(v'\in X^MN_c\). Saturation gives a difference of corrections in
\(DP_c\). This also proves compatibility with maps and addition
of the congruences, by applying the equations before comparing
corrections. For a change of the cut vector by \(u\in XN_c\) the
correction to \((\gamma-1)u\) can be taken to be
\(\ell^\sigma\Phi^j u\), since one may use \(v=u\), \(K=0\)
in those congruences. This is the comparison modulo \(XN_c\)
needed in the Herr assembly.

\end{proof}
\paragraph{The cyclotomic level models.}\label{par:wach-level-models}
Fix the same crystalline lattice map \(f:T_1\to T_2\) as in
Section~\ref{sec:inert-long-orbit}. Write \(K_O=\operatorname{Frac}(O)\)
for the coefficient field, and retain the choice containing the
cyclotomic roots. Consider an arc of length \(a\) from \(c-a\) to
\(c\). At backward offset \(s\), the letter is \(c-s\). A jump of
the decreasing Hodge filtration at \(u\) means that
\(\operatorname{Fil}^u\ne\operatorname{Fil}^{u+1}\).
Let \(G_0\le g\) bound the absolute jumps at thin letters. We
can take \(G_0=g\) for the given thin box.

Write \(b_s=s_{c-s}\)
for the size of a position with larger jumps. Upper and lower bounds \(m_s^+\)
and \(-m_s^-\) for the jumps, and the exponent of the arc
denominator, are
\[
\begin{array}{c|ccc}
 & m_s^+ & m_s^- & H_s^{\mathrm{arc}}=g+m_s^+\\ \hline
 \text{thin} &G_0&G_0&g+G_0\\
 \text{up}   &b_s+S_0&\eta_0&g+b_s+S_0\\
 \text{down} &\eta_0&b_s+S_0&g+\eta_0 .
\end{array}
\]
Here up means that the large jumps are in \([b_s,b_s+S_0]\).
For down they are in \([-b_s-S_0,-b_s]\).
The arc exponents are used in
\[
 \Lambda=\prod_{s<a}p_s^{H_s^{\mathrm{arc}}},\qquad
 \Theta=\prod_{s<a}p_s^g .
\]
The symmetric widths \(H_i\) from the recurrence for \(R\) control
the ordinary Wach cokernel.

For the evaluation at offset \(s\), use the product over all
primitive \(\ell^{s+1}\)-th roots \(\zeta\) of the space of Laurent
series in \(t\) with coefficients in the terminal crystalline
space \(D_{i,c}\). In each component expand the ordinary
cyclotomic parameter at \(1+\pi=\zeta\exp(t)\).
The scalar evaluation map is denoted by \(\rho_s\).

At every primitive root indexing \(\rho_s\) there is an expansion
\[
 \rho_s(p_s)=t\,u_s(t),\qquad u_s(0)\ne0.
\]
Indeed the primitive roots are simple
roots of the cyclotomic polynomial in characteristic zero and
\(\zeta\exp(t)\) has a nonzero derivative at zero.
The tame parameter \(X\) and all \(p_{s'}\) with \(s'\ne s\)
evaluate to units. Frobenius on the series sends
\(\zeta\) to \(\zeta^\ell\) and rescales \(t\) by \(\ell\),
and on the crystalline coefficients it transports them to the
next letter. Cyclotomic action on the degree \(k\) coefficients
is the permutation of the roots given by the cyclotomic value
times its \(k\)-th power. The level space consists of the tuples
fixed by the tame group for this action.

If \(\operatorname{Fil}^u_{c-s}\) denotes the filtration at the starting letter, the filtration at offset \(s\) is
\[
 \mathcal F_{i,s}^u=\Phi_{\rm cris}^s(\operatorname{Fil}^u_{c-s})\subseteq D_{i,c}.
\]
This transport is the de Rham comparison at that level. We label a Wach
step by its terminal Hodge letter. To translate
the de Rham convention using \(\varphi^{-(s+1)}\), label a
Wach component by the embedding one step before its
ordinary coefficient embedding and transport its crystalline
coefficients by \(\Phi_{\rm cris}^{-1}\). An evaluation at
level \(s=0\) then has the Hodge filtration of its own letter.
At level \(s\) the remaining \(s\) steps give the filtration of
\(c-s\), transported by \(\Phi_{\rm cris}^s\).
This convention is common to the source, target and extension.
The Wach evaluation generates the completed rational module
of tuples having coefficient of degree \(k\) in
\(\mathcal F_{i,s}^{-k}\). We denote this condition by
\(J_{i,s}\).

\begin{lemma}[The completed Wach lattice at a cyclotomic root]\label{lem:wach-evaluated-hodge}
Label a Wach step by its terminal Hodge letter and use
\(\Phi_{\rm cris}^{s}\) to transport the filtration from \(c-s\).
If \(\widehat N_{i,s}^{\rm rat}\) denotes the completed rational Wach lattice at cyclotomic level \(s\), then
\[
 \widehat N_{i,s}^{\rm rat}=J_{i,s}\qquad(0\le s<a),
\]
with \(J_{i,s}\) as in Paragraph~\ref{par:wach-level-models}. At the starting letter
reduction at \(X=0\), after inverting \(\ell\), gives the crystalline
space. Both comparisons commute with the coefficient-field action.
\end{lemma}

\begin{proof}
First use a common Tate twist to put the
crystalline weights in the positive range of Berger's Wach
comparison \cite[Propositions II.2.1, II.3.1 and III.2.1]{berger-wach}.
Over the ordinary cyclotomic parameter, the comparison at a
root is the de Rham embedding indexed by that root, with
\(t\) rescaled and crystalline coefficients transported by
Frobenius. If \(d_j\) is an adapted basis with jumps \(u_j\), the completed de Rham lattice is
\[
 \bigoplus_j K_O[[t]]\,t^{-u_j}d_j
\]
in a root component with coefficient field \(K_O\). Indeed these elements lie in the de Rham
lattice and their reductions form a basis by the Hodge--Tate
comparison for the de Rham representation. The completed
Wach module maps into this lattice through
\(N(V)\subset D^+(V)\) in the positive range. Relative to the
lattice generated by the constant crystalline basis, its
elementary-divisor exponents at the root are \(u_j\) by
Berger's two elementary-divisor propositions.

Let \(L\subseteq L'\) denote the inclusion at the root. Their determinant orders relative to the same constant basis are both \(\sum_j u_j\), so
\[
 \operatorname{length}(L'/L)
  =\operatorname{ord}\det L-\operatorname{ord}\det L'=0.
\] This comparison
can first be made on the underlying \(\Q_\ell\)-representation.
It commutes with the coefficient-field action and hence with
the idempotents splitting the coefficients and the embeddings.

Undoing the Tate twist shifts the jumps and the Laurent orders
by the same amount, since the change of Wach basis and of
crystalline basis differ by a power of \(t/\pi\). At a primitive root \(\zeta\ne1\) the local parameter \(t\) vanishes and \(\pi(0)=\zeta-1\ne0\). Hence
\[
 \operatorname{ord}_t(t/\pi)=1.
\]
This is the shift in Laurent order when the Tate twist is undone.

Passage to tame invariants
uses the averaging idempotent for the group of order
\(\ell-1\). Averaging the ordinary basis preserves its reduction
modulo \(\pi\), where the action is trivial. The averaged basis
therefore remains a basis over \(O[[\pi]]\), and its invariant
coefficients are in \(\mathcal R\). This yields the asserted
coefficient condition on the tame-invariant spaces. At the root one, the starting-letter identification is
\[
 N_{i,c-a}[1/\ell]/XN_{i,c-a}[1/\ell]\simeq D_{i,c-a}.
\]
It is the crystalline comparison at the origin
\cite[Corollary III.4.5]{berger-wach}. In crystalline
coordinates the ordinary expansion at one is regular, and
the factor \(t/\pi\) for the twist is a unit there.
\end{proof}

\begin{proposition}[An integral membership test along an arc]\label{prop:wach-arc-membership}
Write \(w_s\) for the level-\(s\) evaluation and put
\[
 W_i=\Lambda^{-1}\operatorname{span}_{\mathcal R}\Phi^a N_{i,c-a}.
\]
Using \(m_s^\pm\) and \(\Lambda\) from
Paragraph~\ref{par:wach-level-models}, one has
\[
 \Bigl(\prod_{s<a}p_s^{m_s^+}\Bigr)N_{i,c}
    \subset\operatorname{span}_{\mathcal R}\Phi^aN_{i,c-a},
 \qquad
 w\in N_{i,c}\ \Longleftrightarrow\
 (\forall s<a)\ w_s\in J_{i,s}\quad(w\in W_i).
\]
The signed step bound applies also to the crystalline extension of
\(O\) by \(T_2\).
\end{proposition}

\begin{proof}
\proofstep{Clearing the one-step denominators.}
These comparisons give integral step and level tests. For
one step ending at \(i\), the span \(\mathcal R\Phi N_{i-1}\)
evaluates to the regular lattice at the first cyclotomic
level. This follows by transport of the comparison at one.
The condition at that level on \(N_i\) is the Hodge filtration
at the ending letter. At the first level the two factors move in opposite directions:
\[
 q^{m^+}N_i\ \subseteq\ \mathcal R\Phi N_{i-1},\qquad q^{m^-}\mathcal R\Phi N_{i-1}\ \subseteq\ N_i,
\]
initially in the evaluated modules. For integral inclusions use the coarse Wach height bounds.
After a common twist the ordinary positive Wach module is
stable under Frobenius, with linearized cokernel killed by
the cyclotomic factor to the power of the width of the
jump box \cite[Definition III.4.1 and Proposition III.4.2]{berger-wach}.
Undoing the twist multiplies the Frobenius matrix by a
power of that factor since \(\phi(\pi)/\pi\) is the ordinary
factor.

Tame averaging gives the corresponding inclusions
on invariant modules. Thus the coordinates on either basis
have an a priori denominator a power of \(q\). The
local test removes that denominator integrally. In fact the
remainder on division by the distinguished associate of
\(q^m\) vanishes when its Taylor coefficients of order
less than \(m\) vanish at the roots. Indeed let \(D\) be the distinguished polynomial and \(R\) the remainder. Over \(K_O\) the Taylor conditions imply
\[
 D\mid R,\qquad \deg R<\deg D\quad\Longrightarrow\quad R=0.
\]
The quotient by \(D\) is free over \(O\), so this vanishing descends integrally. Iterating the resulting signed bound gives
\[
 \Bigl(\prod_{s<a}p_s^{m_s^+}\Bigr)N_{i,c}
    \subset\operatorname{span}_{\mathcal R}\Phi^aN_{i,c-a}.
\]
The same proof applies to the crystalline extension
\(0\to T_2\to T'\to O\to0\). Exactness of the crystalline filtration
adds only the jump at zero, which lies in each of the stated boxes.

\proofstep{Testing the product of the denominators.}
Put
\[
 W_i=\Lambda^{-1}\operatorname{span}_{\mathcal R} \Phi^aN_{i,c-a}.
\]
Using the opposite step bound, there is a product
\(Q_*\) of powers of the factors \(p_s\) with \(Q_*W_i\subset N_{i,c}\).
For \(w\in W_i\), expand \(Q_*w\) on a Wach basis at \(c\).
At each level the Wach basis is a basis of the completed Hodge
lattice at each root. Thus the evaluation of \(w\) is in \(J_{i,s}\)
exactly when its multiplied coordinates for \(Q_*w\) vanish to the
order of \(\rho_s(Q_*)\) at every root of level \(s\).

For the scalar
remainder test over \(O[[X]]\), the values of \(X\) at primitive
roots identify the distinct tame orbits, with one simple root of the
distinguished associate of \(p_s\) per orbit. Weierstrass division
and the same root test therefore give divisibility by the entire
product \(Q_*\) in every coordinate. Consequently
\[
       w\in N_{i,c}\quad\Longleftrightarrow\quad
       w_s\in J_{i,s}\text{ for every }s<a .
\]
The units at each level give the two tests used below. The inverse
image of \(N_{i,c}\) under \(\Theta\) lies in \(W_i\) by the signed
step bound. In \(W_i\) its test is \(p_s^g w_s\in J_{i,s}\)
for every \(s\). Also \(XN_{i,c}\) consists of the vectors of
\(XW_i\) having evaluations in every \(J_{i,s}\); apply the
previous test after dividing by \(X\), a unit at the evaluated
levels.
\end{proof}

Define \(D_{i,s}\) to be the diagonal crystalline constants.
For an up letter write \(U^K_{i,s}=\mathcal F_{i,s}^{b_s}\),
and for down write \(U^K_{i,s}=\mathcal F_{i,s}^{-\eta_0}\).
Define \(B_{i,s}\) by the following increasing degree filtration
on its coefficients. In a thin letter it is zero below degree
zero and full from zero. In an up letter the breaks are \(-b_s\)
and zero. In a down letter they are zero and \(b_s\). In the
last two cases it is zero before the first break, \(U^K_{i,s}\)
from the first to the second break, and full from the second.
Each coefficient condition is imposed at all the roots.

\begin{lemma}[Comparison with the rounded degree filtration]\label{lem:wach-rounded-breaks}
For the coefficient spaces \(B_{i,s}\) with the rounded breaks
just defined, one has
\[
 p_s^gB_{i,s}\subset J_{i,s},\qquad p_s^gJ_{i,s}\subset B_{i,s}.
\]
In coordinates the rounded filtration gives the divisibility conditions
of \eqref{eq:wach-level-condition}, compatibly with the crystalline map.
\end{lemma}

\begin{proof}
Between the two jump windows the Hodge filtration is constant,
with middle subspace precisely \(U^K_{i,s}\). Its two groups
of degrees differ from the rounded breaks by at most \(g\)
since \(G_0,S_0,\eta_0\le g\). In the thin case the one
group differs from zero by at most \(G_0\). It follows,
using multiplication by the unit part of \(\rho_s(p_s)\) on
the increasing filtrations, that
\[
       p_s^g B_{i,s}\subset J_{i,s},\qquad
       p_s^g J_{i,s}\subset B_{i,s}.
\]

Fix a Wach basis \(e_i\) at the start and write \(\lambda_i\)
for its reduction at one transported by \(\Phi^a\). This is
a basis of \(D_{i,c}\). Viewing the basis as a row, the coordinate evaluation at offset \(s\) is
\[
 y\longmapsto\rho_s(\Lambda)^{-1}\lambda_i\rho_s(y),\qquad y\in\mathcal R^{r_i}.
\]
The saturated inverse image \(U_{i,s}\) under \(\lambda_i\) has an integral complement \(C_{i,s}\). Thus
\[
 O^{r_i}=U_{i,s}\oplus C_{i,s},\qquad U_{i,s}=\lambda_i^{-1}(U^K_{i,s})\cap O^{r_i}.
\]
Multiplying
the rounded conditions by \(\rho_s(\Lambda)\) shifts the breaks
by \(H_s^{\mathrm{arc}}\). In bases of the saturated submodule
and its complement, Weierstrass division gives exactly the
three divisibility tests in \eqref{eq:wach-level-condition}.

For a constant plus this lattice, subtract the series
\(\Lambda d\) for the corresponding crystalline constant
in \(K_O^{r_i}\) and clear its finitely many coefficient
denominators. This gives the test with \(\varpi^k y-\Lambda d\)
in that paragraph. The crystalline map on the transported
bases has the matrix of the Wach map reduced modulo \(X\), because
reduction, crystalline comparison and Frobenius are functorial.
Thus both evaluation maps commute with the specified maps.
For the crystalline map \(F\), compatibility of the up submodules comes from
\[
 F(\mathcal F_{1,s}^{b_s})\subseteq\mathcal F_{2,s}^{b_s}.
\]

\end{proof}

\paragraph{One integral partner at all levels of an arc.}\label{par:wach-partners}
The truncation modulus after the arc is
\[
 D'=\phi^{a-j}\left(X^M\prod_{s<j}p_s^{3g}\right).
\]
Its orders at the evaluated levels are
\[
 B_s^*=\begin{cases}
 (\ell-1)M,&s<a-j,\\
 3g,&a-j\le s<a .
 \end{cases}
\]

\begin{lemma}[One integral representative for all levels of an arc]\label{lem:wach-arc-partners}
For the cut and arc of the construction, assume \((\ell-1)M>h+2g\),
\(3g\le(\ell-1)M\), and the stated buffer of thin letters.
The comparison uses the filtration
\[
 F_{i,s}=\{z:\operatorname{ord}(z)\ge B_s^*-H_s^{\rm arc}\}.
\] After the scalar loss \(\ell^\sigma\)
of the cut congruences, a vector in either the Wach or the coordinate
description has one integral partner in the other description
agreeing modulo \(F_{i,s}\) at every \(s<a\).
\end{lemma}

\begin{proof}
\proofstep{Orders of the truncation factors.}
The polynomial \(\phi^{a-j}(X)\) has a simple \(X\)-factor and
\(p_s\)-orders \(\ell-1\) for \(s<a-j\). Define \(F_{i,s}\)
to consist of tuples of order at least \(B_s^*-H_s^{\mathrm{arc}}\).
At each of the last \(j\) backward offsets the letter is thin
by the buffer after the starting cut. At thin offsets the order
is at least \(g\), and at up offsets it is at least \(g+1\).

At down offsets it is at least \(b_s+S_0+1\).
Here we use
\[
 (\ell-1)M>h+2g
\]
and \(b_s+S_0\le h\).
All these degrees are in the full-coefficient part of
both filtrations. Thus \(F_{i,s}\) is contained in
\(J_{i,s}\cap B_{i,s}\). The same order is used for the two
representations and the crystalline map preserves it.
If a coordinate evaluation is in \(F_{i,s}\), its
numerator is divisible by \(p_s^{B_s^*}\) by the scalar
remainder test, in particular by \(p_s^{H_s^{\mathrm{arc}}}\).

\proofstep{A basis fixed modulo the cut.}
There is a basis change \(A_i(X)\) modulo \(X^M\), with
constant term the identity, making \(e_i A_i\)
\(\gamma\)-invariant modulo \(X^M\) over \(K_O\).
Construct it degree by degree, dividing at degree \(r\) by
\(\omega_X^r-1\), just as in the cut correction.
Write
\[
 A_i(X)=1+\sum_{r=1}^{M-1}A_{i,r}X^r.
\]

If \(C_{i,r}\) is the degree-\(r\) coefficient of the remaining
invariance defect of the previous truncation in the basis \(e_i\),
the next correction solves
\[
  (\omega_X^r-1)A_{i,r}=-C_{i,r},\qquad
  A_{i,r}\in\ell^{-a_r}\operatorname{Mat}(O),\qquad
  a_r=\sum_{k=1}^r(1+v_\ell(k)).
\]
For the inverse, every term in degree \(r\) of the geometric series
has the form \(\pm A_{i,r_1}\cdots A_{i,r_t}\), where
\(r_1+\cdots+r_t=r\). For a product of coefficients with degrees \(r_1+\cdots+r_t=r\), the extra valuation is controlled by
\[
 \sum_j v_\ell(r_j!)\le v_\ell(r!),\qquad \frac{r!}{\prod_j r_j!}\in\mathbb Z.
\]

In subsequent bounds write \(A_r=a_r\). Use polynomial representatives of the two
truncated matrices. At the root one the ordinary parameter
satisfies \(X\sim t^{\ell-1}\). In crystalline coordinates the
normalized basis at one is constant modulo \(t^{(\ell-1)M}\):
if \(u_k\) is a coefficient in an intermediate degree, cyclotomic invariance gives
\[
 (\omega^k-1)u_k=0,\qquad 0<k<(\ell-1)M ,
\]
and \(\omega^k\ne1\) forces \(u_k=0\). Every evaluated root reaches one after
\(a\) Frobenius steps, which transport this constant to
\(\lambda_i\) and preserve the orders of the nonconstant terms.

\proofstep{The two integral representatives.}
Applying \(\phi^{a-j}\) to the scalar cut congruences shows that
\(\ell^\sigma\phi^a(A_i)\) and
\(\ell^\sigma\phi^a(A_i^{-1})\) are congruent to integral
matrices modulo \(D'\) over \(\mathcal R[1/\ell]\).
Here each coefficient uses the bound with \(A_r\) at its own
degree. Also \(\phi^a(X)^M\) is divisible by \(D'\)
by the upper bound
\[
 3g\le(\ell-1)M.
\]
Consequently
the product of the two transported matrices is the identity
modulo \(D'\) before multiplying by the scalar.
Write an arc vector as \(\Lambda^{-1}(\Phi^a e_i)v\).
Its coordinate partner is obtained by applying an integral
representative of \(\ell^\sigma\phi^a(A_i^{-1})\) to \(v\).

Conversely, the arc partner of a coordinate vector \(y\)
uses \(v\) obtained by applying an integral representative
of \(\ell^\sigma\phi^a(A_i)\) to \(y\).
In either direction the transported invariant basis gives
the required agreement after multiplying by \(\Lambda\), modulo
\(t^{B_s^*}\) at every level. Dividing by \(\Lambda\) gives
the congruence modulo \(F_{i,s}\). Each partner uses a single
integral vector for the whole arc.
\end{proof}

\paragraph{The two cyclotomic polynomial actions.}\label{par:wach-polynomial-actions}
Let \(\omega=\chi(\gamma)\) and let \(\tau_s\) be the permutation
of the primitive roots at level \(s\) given by \(\omega\).
On degree \(k\), \(\gamma\) acts by \(\omega^k\tau_s\).
The principal units act through a cyclic group of order
\(\ell^s\), generated by \(\tau_s\), and together with
the tame units act transitively on the roots.

\begin{lemma}[Polynomial corrections at all evaluated levels]\label{lem:wach-level-polynomials}
On the tame-invariant space,
\[
 (\gamma-1)w_s\in J_{i,s}\quad\Longrightarrow\quad w_s\in D_{i,s}+J_{i,s}.
\] For \(E=a-1\) the polynomials
\(P_{E,I}\) constructed from the respective target and source
degree intervals below send \(p_s^{-g}B_{2,s}\) into
\(D_{2,s}+B_{2,s}\) and \(p_s^{-g}J_{1,s}\) into \(D_{1,s}+J_{1,s}\)
at all offsets. The degrees used lie in at most three intervals
of length \(3g\) inside the prescribed range \([-2h,2h]\).
\end{lemma}

\begin{proof}
Suppose \((\gamma-1)w_s\in J_{i,s}\). In degree \(k\ne0\),
let \(\bar w_{s,k}\) be the fixed class in the degree-\(k\)
quotient. The permutation of the roots has order \(\ell^s\), so
\[
   \bar w_{s,k}=\gamma^{\ell^s}\bar w_{s,k}
               =\omega^{k\ell^s}\bar w_{s,k}.
\]
For \(k\ne0\), \(\omega^{k\ell^s}\ne1\), hence
\(\bar w_{s,k}=0\) over the coefficient field. In degree zero the class is both tame and
principal invariant, thus constant across the roots. A lift
of that constant in \(D_{i,c}\) shows \(w_s\in D_{i,s}+J_{i,s}\).
The action fixes diagonal constants and preserves \(J_{i,s}\).

\proofstep{Averaging the degree-zero coefficient.}
Use the polynomials \(P_{E,I}\) of the arc calculation with
\(E=a-1\), so \(E\ge s\) at each level.
A factor \(Z^{\ell^E}-\omega^{k\ell^E}\) kills degree \(k\).
In degree zero the sum \(\sum_{v=0}^{\ell^E-1}Z^v\) averages
over the principal orbit with a scalar multiplicity.
For a tame-invariant degree-zero tuple \(b\) on the roots, its
\(\gamma\)-orbit has period dividing \(\ell^s\). Hence
\[
 \sum_{v=0}^{\ell^E-1}\gamma^v b
   =\ell^{E-s}\sum_{j=0}^{\ell^s-1}\gamma^j b
   \qquad(E\ge s).
\]
The sum is fixed by \(\gamma\) and by the tame action. Their transitivity
makes it diagonal constant, a space preserved by the remaining factors.

\proofstep{The rounded target and the filtered source.}
For the target polynomial, write \(\mathcal B_s\) for the rounded breaks defining \(B_{2,s}\). Take precisely the integer degrees
\[
 I=\bigcup_{e\in\mathcal B_s}\bigl([e-g,e)\cap\mathbb Z\bigr).
\]

A vector in \(p_s^{-g}B_{2,s}\) satisfies in degree \(k\)
the rounded condition for \(k+g\).
When this condition is larger than that for \(k\), one of
the selected intervals contains \(k\). Every such \(k\)
is nonzero, since \(b_s>g\) on positions with larger jumps.
The polynomial kills these degrees, and in the other degrees
it preserves the required subspace. It fixes the space of
constants up to scalar, giving the required map to
\(D_{2,s}+B_{2,s}\).

On the source side use the separate set
\[
 I_{\rm src}=\bigcup_{u}\bigl([-u-g,-u)\cap\mathbb Z\bigr),
\]
where \(u\) runs over the jumps in the prescribed source boxes.
The degree test for \(p_s^{-g}J_{1,s}\) compares
\(\mathcal F_{1,s}^{-k-g}\) with \(\mathcal F_{1,s}^{-k}\).
Every change occurs in a selected interval. The factors kill
all offending nonzero degrees. In degree zero the principal
sum gives a diagonal constant, which is allowed in
\(D_{1,s}+J_{1,s}\). At degrees where the subspace is
unchanged it is preserved. This proves the second action
on the same evaluation model.

For each polynomial the union lies in three integer intervals
of length at most \(3g\). After enlarging \(G_0\) to include \(\eta_0\), the central interval and its length satisfy
\[
 I_{\rm cent}\subseteq[-G_0-g,G_0-1],\qquad |I_{\rm cent}|\le2G_0+g\le3g.
\]
The intervals around a
thick up or down break lie within the corresponding one-sided
translated window of length at most \(g+S_0\).
The arc contains at most two positions with larger jumps. All degrees lie
in \([-2h,2h]\). These are precisely the window and order
bounds used for the scalar valuation at one. The powers of
\(\ell\) in the partners and polynomial values use only the
maximum length of an arc and the prescribed boxes.
\end{proof}


\paragraph{A level comparison for the mixed-sign arc.}
\label{par:wach-mixed-constant}
Use the arc length \(a\le A\), the powers \(H_s^{\mathrm{arc}}\),
and the level data of Paragraph~\ref{par:wach-level-models}.
Distinguish the transported ranks \(r_i\), the coordinate map
\(L:O^{r_1}\to O^{r_2}\) on the starting bases modulo \(X\),
and its coefficientwise extension \(L_{\mathcal R}\).
The polynomial comparison uses, at each \(s\), evaluation
maps from \(W_i\) and from \(\mathcal R^{r_i}\) to a common level
space. Denote their values by \(w_s\) and \(y_s\), respectively.
The latter refers to a vector \(y\) of coordinate series, not to
its coefficients at \(X=0\). The following are the lattice
properties used in this comparison.

\begin{itemize}
\item There are two local lattices \(B_{i,s},J_{i,s}\) with
\(p_s^g B_{i,s}\subset J_{i,s}\) and
\[
 p_s^g J_{i,s}\subset B_{i,s}.
\]
A submodule \(F_{i,s}\) lies
in both. Every vector on either side of the evaluation has a partner
on the other side, whose evaluation agrees with \(\ell^\sigma\)
times the original modulo \(F_{i,s}\), simultaneously at all
levels. The two evaluation maps commute with the corresponding
maps induced by \(f\) and \(L_{\mathcal R}\). The level map
preserves the \(F_{i,s}\).
\item The preimage of \(N_{i,c}\) under multiplication by
\(\Theta\) lies in \(W_i\). On \(W_i\) this preimage is tested
by
\[
 p_s^g w_s\in J_{i,s}
\]
for all \(s\).
The map \(f\) lifts on this preimage with cost \(\ell^R\).
If \(y_s\in F_{2,s}\), then \(p_s^{H_s^{\mathrm{arc}}}\) divides
\(\varpi^k y\) for some \(k\), where \(\varpi\) is a
uniformizer of \(O\).
\item There are saturated constant submodules \(U_{i,s}\) of
\(O^{r_i}\) and subgroups \(D_{i,s}\) of crystalline constants
in the level spaces. Membership \(y_s\in B_{i,s}\) is equivalent,
over all levels, to the conditions
\[
 y\in
 \begin{cases}
 p_s^{H_s^{\mathrm{arc}}}\mathcal R^{r_i} &\text{smaller jumps},\\
 p_s^{H_s^{\mathrm{arc}}-b_s}\mathcal R U_{i,s}+p_s^{H_s^{\mathrm{arc}}}\mathcal R^{r_i}
                          &\text{up},\\
 p_s^{H_s^{\mathrm{arc}}}\mathcal R U_{i,s}+p_s^{H_s^{\mathrm{arc}}+b_s}\mathcal R^{r_i}
                          &\text{down},
 \end{cases}
 \tag{\(*\)}\label{eq:wach-level-condition}
\]
after allowing multiplication of \(y\) by a power of \(\varpi\)
at each level. Membership \(y_s\in D_{i,s}+B_{i,s}\) is tested
by the same conditions on \(\varpi^k y-\Lambda d\) at each level,
for some \(k\ge0\), \(d\in O^{r_i}\). Write \(\mathcal N_i\)
and \(\mathcal I_i\) for these two modules of coordinate series.
\end{itemize}

\begin{lemma}[The up subspace after a down letter]\label{lem:up-after-down}
Suppose the two positions with larger jumps are a down letter
at \(s_0\) and an up letter at \(s_1>s_0\).
Assume \(L(U_{1,s_1})\subset U_{2,s_1}\),
\(2b_{s_0}\le\ell b_{s_1}\) and \(2a_0<b_{s_1}\), where
\(a_0=2\sigma+R+2gA\) is the initial exponent of
Section~\ref{sec:inert-long-orbit} and \(a\le A\) is the arc length. Then
\[
 \ell^{2a_0}U_{2,s_1}\subset L(U_{1,s_1}).
\]
\end{lemma}

\begin{proof}
Put \(Q=p_{s_1}^{b_{s_1}}\), \(P=p_{s_0}^{b_{s_0}}\) and
\(\Lambda_1=\Lambda/Q\). If \(v\in U_{2,s_1}\), the series
\(y_2=P\Lambda_1 v\) belongs to \(\mathcal N_2\), directly
by \eqref{eq:wach-level-condition} at each of the two positions with larger jumps and at every position with smaller jumps.

Choose its partner \(w_2\in W_2\). The error at each level lies
in \(F_{2,s}\) and \((y_2)_s\in B_{2,s}\), so
\[
 p_s^g(w_2)_s\in J_{2,s}.
\]
The preimage test gives a lift
\(w_1\) on the \(\Theta\)-divisible lattice, with
\(f(w_1)=\ell^R w_2\). Choose the coordinate partner \(y_1\)
of \(w_1\). Each \(p_s^{2g}(y_1)_s\) belongs to \(B_{1,s}\),
so \(\Theta^2 y_1\in\mathcal N_1\).
The partner errors and compatibility with \(f\) give
\[
 (L_{\mathcal R}y_1-\ell^{2\sigma+R}y_2)_s\in F_{2,s}.
\]

Every \(p_s^{H_s^{\mathrm{arc}}}\) therefore divides the difference after
multiplication by a power of \(\varpi\). These divisibilities
hold integrally: each quotient by a distinguished polynomial
power is free over \(O\) by Weierstrass division. Their product
\(\Lambda\) divides by relative primality after inverting
\(\varpi\), followed by the same saturation for the product.
This uses no integral Chinese remainder decomposition.

\proofstep{Testing cyclic quotients of the image.}
Take an integral linear form \(\lambda\) on \(O^{r_2}\)
such that \(\lambda(v)=1\) and
\(\lambda L(U_{1,s_1})\subset tO\).
For any \(y\in\mathcal N_1\),
\[
                 \lambda L_{\mathcal R}(y)
                   \in \Lambda_1(t,Q).
\]
Indeed split \(O^{r_1}\) along the saturated \(U_{1,s_1}\).
On its complement the up condition implies divisibility by
\(p_{s_1}^{H_{s_1}}\), and on the summand \(U_{1,s_1}\) it
implies divisibility by \(p_{s_1}^{H_{s_1}-b_{s_1}}\).
All coordinates are divisible by \(p_s^{H_s^{\mathrm{arc}}}\) at the other
levels. Multiplying the powers of the distinct cyclotomic factors, the complementary
summand gives a multiple of \(\Lambda\) and the up summand gives
a multiple of \(t\Lambda_1\). Applying this to
\(\Theta^2 y_1\) and using the congruence for
\(L_{\mathcal R}y_1\) yields, after cancelling \(\Lambda_1\),
\begin{equation}\label{eq:wach-gauss-ideal}
                 \ell^{2\sigma+R}\Theta^2P\in(t,Q).
\end{equation}

\proofstep{The Gauss valuation and Smith factors.}
Normalize \(v_\varpi(\varpi)=1\), and put \(e_O=v_\varpi(\ell)\).
For \(x>0\) the Gauss valuation of \(\sum a_m X^m\) is
\(\min_m(v_\varpi(a_m)+mx)\). If \(i,j\) are the least minimizing indices for \(f,g\), then in
the coefficient of \(X^{i+j}\) only \(f_i g_j\) attains the minimum.
Indeed a different pair with the same index sum uses an index below
\(i\) or below \(j\). Therefore
\[
 v_\varpi([X^{i+j}]fg)+(i+j)x=v_x(f)+v_x(g)=v_x(fg).
\] Preparation gives
\(v_x(p_s)=\min(e_O,\ell^s x)\).

If \(t\ne0\), set
\(k=v_\varpi(t)\); for \(t=0\) take \(k=+\infty\).
The assertion for \(k=0\) is immediate. Otherwise set
\(b_1=b_{s_1}\), \(b_0=b_{s_0}\),
\(m=\min(k,e_O b_1)\), and \(x=m/(b_1\ell^{s_1})\). Then
\(v_x(Q)=m\), and \(v_x(P)\le b_0 m/(\ell b_1)\).
The ideal relation and \(v_x(p_s)\le e_O\) give
\[
 m\le e_O(2\sigma+R+2ga)+\frac{b_0}{\ell b_1}m
   \le e_O a_0+m/2 .
\]
Thus \(m\le2e_Oa_0<e_Ob_1\), which excludes \(t=0\) and
proves \(t\mid\ell^{2a_0}\).

Since \(U_{2,s_1}\) is a direct
summand, Smith coordinate forms extend to the ambient
\(O^{r_2}\). Apply the divisibility assertion to its Smith
basis vectors for the inclusion of \(L(U_{1,s_1})\).
A missing image direction would give \(t=0\); the invariant
factors in the image directions all divide \(\ell^{2a_0}\).
Consequently
\[
                   \ell^{2a_0}U_{2,s_1}\subset L(U_{1,s_1}).
\]
The use of \(s_0<s_1\) gives the factor \(1/\ell\) in the
valuation comparison. The argument measures divisibility by a
scalar in a ramified coefficient ring and so is uniform under
the coefficient extension.
\end{proof}

\paragraph{From constant coordinates to the arc lattice.}
\begin{lemma}[Lifting constant coordinates across an arc]\label{lem:wach-constant-lift}
With the integral lattices and scalar \(a_0\) of this comparison,
each \(y\in\mathcal I_2\) has a lift modulo \(\mathcal N_2\):
\begin{equation}\label{eq:wach-const-coord-lift}
 \ell^{2a_0}y=L_{\mathcal R}y_1+\ell^{2a_0}\nu,\qquad
 y\in\mathcal I_2,\quad y_1\in\mathcal I_1,\quad
 \nu\in\mathcal N_2 .
\end{equation}
\end{lemma}
\begin{proof}
Here is the constant comparison for all the words. For either
representation, omit its index and let \(V=O^r\). If there is
no down letter, the level conditions are unchanged by adding a
local constant times \(\Lambda\), so \(\mathcal I=\mathcal N\).
This includes a word with two up letters. Otherwise there is at
most one up letter. Write \(Q=p_s^{b_s}\) at the up letter and
\(Q=1\) if there is none. Divisibility at every level gives
\(y=(\Lambda/Q)Y\), with \(Y\in\mathcal R^r\).
At an up letter the condition for \(Y\) is
\[
             Y\in\mathcal R U_s+Q\mathcal R^r .
\]

At a down letter with \(P=p_s^{b_s}\) it is the condition that
the image of \(Y\) in \((\mathcal R/(P))\otimes_O(V/U_s)\)
lie in the rational constant line factor \(\overline Q\) times
\((V/U_s)\otimes_O\operatorname{Frac}(O)\). To see this,
divide by the full local power \(p_s^{H_s^{\mathrm{arc}}}\), then cancel the
other level factors modulo \(P\) after inverting \(\varpi\).
The difference between \(\mathcal I\) and \(\mathcal N\) is
that \(\mathcal N\) requires the image at each down letter to
be zero. The quotients in this test are free over \(O\) by
Weierstrass division and saturation of \(U_s\).

With a single down letter and no up letter, primitivity means the coefficient map splits the constant inclusion:
\[
 O\xrightarrow{\ a\mapsto a\ }\mathcal R/(P)\xrightarrow{\lambda}O,\qquad \lambda(1)=1.
\]
In fact, if the projected class is \(1\otimes v\), choose an integral
linear form \(\lambda\) on \(\mathcal R/(P)\) with \(\lambda(1)=1\).
Then
\[
 (\lambda\otimes1)(1\otimes v)=v\in V/U_s .
\]
This recovers the constant integrally in the saturated quotient. If the up letter
has smaller offset than the down one, its power \(Q\) reduces
modulo \(\varpi\) to a unit times a monomial of smaller degree
than the distinguished polynomial for \(P\), by the inequality
\(b_{\rm up}<\ell b_{\rm down}\). The class of \(Q\) modulo
\(P\) is again primitive, which proves integrality of the
constant by the same linear-form test. In both cases lift the
constant to \(d\in V\). Then
\[
 y-\Lambda d\in\mathcal N.
\]


\proofstep{Two down letters.}
For two down letters let \(P,Q\) be their powers in increasing
offset order, and \(U_0,U_1\) the corresponding submodules.
There are integral constants \(d_0,d_1\in V\) by the separate
primitive test for \(1\) in each quotient. Their difference
belongs to the ideal \((P,Q)\) in
\(\mathcal R\otimes_O V/(U_0+U_1)\). Let \(\rho\) be least with
\[
 \varpi^\rho\in(P,Q),
\]
which exists by coprimality after
inverting \(\varpi\). The constant components in a cyclic
decomposition of \(V/(U_0+U_1)\) belong to
\(\varpi^\rho V/(U_0+U_1)\). On a free summand this follows
from the definition of \(\rho\).

On a summand of order
\(\varpi^k\), if an integral representative \(a\) in
\((P,Q,\varpi^k)\) had valuation below \(\min(k,\rho)\), then
\(a-\varpi^k f\in(P,Q)\) for some \(f\in\mathcal R\).
The factor \(1-(\varpi^k/a)f\) is a unit, a contradiction
to the choice of \(\rho\).
Write
\[
 d_0-d_1=u_0+u_1+\varpi^\rho d'
\]
with
\(u_i\in U_i\). Choose \(f_\rho\in(Q)\) with
\(f_\rho-\varpi^\rho\in(P)\). Put \(Y'=Y-d_1-u_1-f_\rho d'\).
The two reductions are
\[
 \begin{aligned}
  Y'&\equiv d_0-d_1-u_1-\varpi^\rho d'=u_0
        &&\pmod {P\mathcal R^r+\mathcal R U_0},\\
  Y'&\equiv d_1-d_1-u_1=-u_1
        &&\pmod {Q\mathcal R^r+\mathcal R U_1}.
 \end{aligned}
\]
Both images vanish. Multiplying back by \(\Lambda\) gives the
two factors \(\Lambda\) and \(\Lambda f_\rho\).

\proofstep{A down letter before an up letter.}
For a down letter at the smaller offset and an up letter at
the larger one, write \(U=U_{s_1}\) and \(V_0=U_{s_0}\),
with \(P,Q,\Lambda_1\) as in the mixed-sign estimate. Write
\(Q=\varpi^k\alpha+P\beta\), with \(k\) largest possible.
There is such a finite \(k\) because the image of \(Q\) in
the finite free \(\mathcal R/(P)\) is nonzero. The class
of \(\alpha\) is primitive. The down condition for \(Y\)
therefore gives
\[
 \overline Y=\overline\alpha d_0
\]
in
\((\mathcal R/(P))\otimes_O(V/V_0)\) for some \(d_0\in V/V_0\).
The up condition, projected further to \(V/(U+V_0)\),
shows that \(\overline\alpha d_0\) belongs to
\(\varpi^k\overline\alpha((\mathcal R/(P))\otimes_O V/(U+V_0))\).

A linear form taking value one on \(\overline\alpha\)
gives \(d_0=v+\varpi^k d\) modulo \(V_0\), with
\(v\in U,d\in V\). This test works on the tensor product
with \(V/(U+V_0)\) even if that module has torsion. Now
\[
 \overline{Y-\alpha v-Qd}
    = \overline\alpha(d_0-v-\varpi^k d)=0
     \quad\text{in }(\mathcal R/(P))\otimes_O(V/V_0).
\]
The same difference still lies in the required up lattice:
\(\alpha v\in\mathcal R U\) and \(Qd\in Q\mathcal R^r\).
Upon multiplying back by the coordinate factors, the resulting
normal form uses \(\Lambda\) times an ordinary constant and
\(\alpha\Lambda_1\) times a vector of \(U\).
In each of the cases the displayed representatives themselves
satisfy the \(\mathcal I\) conditions, by reducing at the same
levels.

On those normal forms, the ordinary constants lift using the
cost \(\ell^R\) on the starting constant map and \(R\le2a_0\).
In the mixed case the additional constant in \(U_{2,s_1}\)
lifts by the Gauss valuation and Smith argument. This gives
\eqref{eq:wach-const-coord-lift}.
\end{proof}

\begin{lemma}[Integral lift between the two polynomial actions]
Let \(x'\) be a target arc vector evaluating in \(D_{2,s}+B_{2,s}\)
at each level of the cut. There is one source vector \(z\) for the arc map \(f\) such that
\[
 f(z)=\ell^{2\sigma+2a_0+R}x',\qquad z_s\in D_{1,s}+p_s^{-g}J_{1,s}\quad\text{at every level }s,
\]
where \(z_s\) is its evaluation.
\end{lemma}
\begin{proof}
Choose the coordinate partner \(y\) of \(x'\). Its error belongs
to \(F_{2,s}\subset B_{2,s}\), so \(y\in\mathcal I_2\).
Apply \eqref{eq:wach-const-coord-lift} to \(y\) and choose the
source arc partner \(x_1\) of \(y_1\). Put
\(\Delta=\ell^{2\sigma+2a_0}x'-f(x_1)\).
In its evaluation the target partner error has coefficient
\(-\ell^{\sigma+2a_0}\), the mapped source partner error has
coefficient \(-1\), and the remaining term is
\(\ell^{\sigma+2a_0}\nu_s\). All belong to \(B_{2,s}\);
the preimage test therefore puts \(\Delta\) in the
\(\Theta\)-divisible lattice. Lift \(\ell^R\Delta\) there to
\(\Delta_1\). Then
\[
 f(\ell^Rx_1+\Delta_1)
 =\ell^R(f(x_1)+\Delta)=\ell^{2\sigma+2a_0+R}x'.
\]
The source partner congruence and
\(B_{1,s}\subset p_s^{-g}J_{1,s}\) give the stated evaluation
at every level, with the same exponent.
\end{proof}
\begin{lemma}[The lift on invariant classes]
The invariant first-coordinate test used in the cut comparison is
\[
 (\gamma-1)z\in XN.
\]
The source and target polynomials in \(\gamma\) and the intervening lattice lift preserve this test in the sense of that comparison.
\end{lemma}
\begin{proof}
To apply this lift, use the level properties of
Paragraphs~\ref{par:wach-level-models} and~\ref{par:wach-polynomial-actions}. The level actions commute with the
evaluation of arc vectors. If the difference under \(\gamma-1\)
is in \(J_{i,s}\) at a level, the vector there lies in
\(D_{i,s}+J_{i,s}\). Writing \(P_2,P_1\) for the target and source polynomials, the actions at a level are
\[
 \begin{aligned}P_2(\gamma)(D_{2,s}+p_s^{-g}B_{2,s})&\subseteq D_{2,s}+B_{2,s},\\ P_1(\gamma)(D_{1,s}+p_s^{-g}J_{1,s})&\subseteq D_{1,s}+J_{1,s}.\end{aligned}
\]
The action fixes \(D_{i,s}\) and preserves
\(J_{i,s}\). Finally \(XN_{i,c}\) consists of the vectors in
\(XW_i\) whose evaluations belong to all the \(J_{i,s}\).
These are statements on the two level models at the same cut.

The action on \(W_i/XW_i\) is trivial. To see this, use the
trivial action on the starting Wach lattice modulo \(X\) and
the commutation with \(\Phi\), together with
\(\phi^a(X)\in(X)\). Division by \(\Lambda\) preserves the trivial action at the origin, since
\[
 \gamma(\Lambda)/\Lambda\equiv1\pmod X.
\]
For an arc vector \(w\) with
\((\gamma-1)w\in XN_{2,c}\), the first level property gives
evaluations in \(D_{2,s}+J_{2,s}\).

Since
\(J_{2,s}\subset p_s^{-g}B_{2,s}\), the target polynomial
produces the input \(x'\) of the intervening lattice lift.
The output, after applying the source polynomial, has
evaluations in \(D_{1,s}+J_{1,s}\). Its \(\gamma-1\) image
therefore has evaluations in all \(J_{1,s}\) and belongs to
\(XW_1\). The last level property puts the image in
\(XN_{1,c}\). An operator polynomial on a vector fixed
modulo \(XN_{2,c}\) agrees with its value at one modulo that
submodule. Writing \(P_2,P_1\) for the target and source
polynomials and \(c_{\rm lift}=2\sigma+2a_0+R\) for the intervening exponent,
choose \(y\) with
\[
 f(y)=\ell^{c_{\rm lift}} P_2(\gamma)w.
\]
Then
\[
 f(P_1(\gamma)y)
 =\ell^{c_{\rm lift}} P_1(\gamma)P_2(\gamma)w
 \equiv\ell^{c_{\rm lift}} P_1(1)P_2(1)w\pmod{XN_{2,c}} .
\]

\end{proof}
\paragraph{The colength at the cuts of the long orbit.}\label{par:wach-cut-colength}

\begin{lemma}[The scalar annihilator at a cut]\label{lem:wach-cut-annihilator}
Let \(f\) be the integral crystalline surjection on the local
arc of Paragraph~\ref{par:wach-level-models}. For \(R\) defined
in Paragraph~\ref{par:inert-cut-exponents}, write \(f_N\) for the Wach map at the cut and \(f_N^{\rm tame}\) for its map on tame invariants. Then
\[
 \ell^R\operatorname{coker} f_N=0,\qquad \ell^R\operatorname{coker} f_N^{\rm tame}=0.
\]
In the extension it gives a lift of \(\ell^R e_0\), and gives
the corresponding constant-map bound after reduction modulo \(X\).
\end{lemma}

\begin{proof}
For an integral crystalline surjection, let \(s_i\) be the
least exponent of the ordinary cyclotomic parameter \(\pi\)
that kills the Wach cokernel at letter \(i\), before tame
invariants. The global height box in this case is symmetric
and has width
\[
 H_{\max}=\max_i H_i=2h.
\]
A common Tate
twist moves both representations into the positive range
of Berger's convention. This multiplies both Wach lattices
by the same power of \(\pi\) in the twisted coordinates and
preserves the cokernel. Berger's integral comparison
\cite[Theorem III.3.1]{berger-wach} places the base change of
each positive Wach lattice between the corresponding base
change of \(T_i\) and its sublattice multiplied by
\(\pi^{H_{\max}}\). Apply the comparison over
\[
 \widetilde{\mathbf A}^{+}=W(\widetilde{\mathbf E}^{+}),
\]
which contains the integral period ring of that comparison.

Since the integral Galois lattice map is onto,
\(\pi^{H_{\max}}N_2\) lies in the image after this base change.
This base change is faithfully flat over
\(\mathcal O_F[[\pi]]\). Modulo \(\ell\) the tilt valuation
ring is torsion-free over the DVR \(k_F[[\overline\pi]]\),
and \(\overline\pi\) is a nonunit in it. The ring of Witt
vectors is \(\ell\)-torsion-free and \(\ell\)-adically complete,
so the adic flatness criterion gives flatness
\cite[Lemma 15.28.5]{stacks-flat}. The residue of the maximal
ideal is proper after extension, giving faithfulness.
Extending by \(O\) and taking the embedding factors preserve
faithful flatness. In the ordinary cokernel
\[
 C_i=N_{2,i}/f(N_{1,i}),
\]
the inclusion
after base change reads
\[
 (\pi^{H_{\max}}C_i)\otimes_{\mathcal O_F[[\pi]]}
      \widetilde{\mathbf A}^+=0.
\]
Faithful flatness implies \(\pi^{H_{\max}}C_i=0\), hence \(s_i\le H_{\max}\). The same argument applies to every crystalline
extension of the trivial lattice by \(T_2\), since its
jumps are the union of the target jumps and zero.

\proofstep{The scalar loss from the cyclotomic action.}
Use a topological generator \(\gamma_*\) of the full
cyclotomic group whose cyclotomic value \(u_*\) generates
\(\Z_\ell^\times\). This action is trivial on Wach lattices
modulo \(\pi\). Suppose \(\ell^k\pi^r N_{2,i}\subset f(N_{1,i})\)
for \(r>0\). For any \(v\in N_{2,i}\), choose \(x\) lifting
\(\ell^k\pi^r v\). The vector \((\gamma_*-1)x/\pi\) is
Wach-integral. Its image, modulo \(\ell^k\pi^rN_{2,i}\), is
\[
                   \ell^k\pi^{r-1}(u_*^r-1)v .
\]
This follows by reducing \(\gamma_*(\pi)/\pi\) and the action
on \(v\) modulo \(\pi\).

The error is already in the image.
One therefore lowers \(r\) by one at the cost
\(v_\ell(u_*^r-1)\). The cost is zero unless
\[
 \ell-1\mid r,
\]
and it is \(1+v_\ell(r/(\ell-1))\)
otherwise. Thus the initial parameter exponent \(s\) gives
a scalar exponent at most \(\ell s/(\ell-1)^2\) by
Legendre's formula.

\proofstep{The recurrence in the letter exponents.}
To pass through Frobenius, put \(d_0=\ell-1\) and let
\(s=s_i\). From \(s\) down to \(r<s\), the exact cost is
\(n-m+v_\ell(n!/m!)\), with \(n=\lfloor s/d_0\rfloor\)
and \(m=\lfloor r/d_0\rfloor\). For \(n=0\) this is zero.
Otherwise it is at most
\[
         \frac{\ell(s-r)}{d_0^2}+\frac{\ell}{d_0}
                         +\lfloor\log_\ell n\rfloor .
\]
Indeed each term in the factorial valuation differs from
\((n-m)/\ell^v\) by less than one, and there are at most
\(\lfloor\log_\ell n\rfloor\) nonzero terms.

Take
\(a=\lceil s/d_0\rceil+c_H+2\). Since
\(q_{\rm cyc}\in(\ell,\pi^{d_0})\), the power
\(\phi(\pi)^a=\pi^a q_{\rm cyc}^a\) is in the ideal
generated by \(\ell^{a-j'}\pi^{a+d_0j'}\) for
\(0\le j'\le a\).
Those terms with \(r=a+d_0j'\ge s\) kill the cokernel.
For \(r<s\),
\[
 a-j'-\frac{\ell(s-r)}{d_0^2}
 =a-\frac{\ell(s-a)}{d_0^2}+\frac{j'}{d_0}
 \ge\frac{a}{d_0^2}+\frac{\ell(c_H+2)}{d_0}.
\]
This exceeds \(\ell/d_0+\lfloor\log_\ell n\rfloor\) since
\(n\le H_{\max}\). Hence each term kills the cokernel by
the cost bound, and so does \(\phi(\pi)^a\).

Write \(H_i=M_i+h_i\) for the letter's width, with the two ordinary one-step height inclusions
\[
 q_{\rm cyc}^{M_i}\Phi N_{2,i-1}\subset N_{2,i},\qquad q_{\rm cyc}^{h_i}N_{1,i}\subset\operatorname{span}(\Phi N_{1,i-1}).
\]

Lift \(\phi(\pi)^a q_{\rm cyc}^{M_i}\Phi v\) for
\(v\in N_{2,i-1}\) and multiply its lift by
\(q_{\rm cyc}^{h_i}\pi^{H_i}\).
It now lies in the source Frobenius span and maps to
\(\Phi(\pi^{a+H_i}v)\).
Frobenius linearization is injective on both lattices,
as it is invertible on the ambient \((\Phi,\gamma)\)-modules
after localization. The scalar ring is free over its
Frobenius image with a basis starting at \(1\), by
Weierstrass division by \(\phi(\pi)-Y\).
Compare the coefficient of \(1\) in the Frobenius-span
expression for the lift.

This constant coefficient gives a lift of
\(\pi^{a+H_i}v\) in the preceding letter. Consequently
\[
                  s_{i-1}\le H_i+c_H+3+\frac{s_i}{\ell-1}.
\]
Iterating the recurrence leaves a remainder bounded by a
geometric fraction of \(H_{\max}\), which tends to zero.
This gives \(s_i\le S(i)\), in addition to
\(s_i\le H_{\max}\). The scalar annihilator above therefore
has exponent at most \(R\) at each cut.

\proofstep{Passage to tame invariants.}
Averaging a lift of an invariant vector by the tame group costs
a unit. Thus the ordinary-module annihilator yields the desired
annihilator on the tame-invariant Wach lattices. This gives
the lift of \(\ell^R e_0\) in the extension Wach lattice and
the ordinary constant-map bound at the preceding cut after
reducing the surjection modulo \(X\). Dividing by \(\Theta\)
commutes with the map and is possible in the ambient module
because the factors of \(\Theta\) are units there. It gives
the required lift on the \(\Theta\)-divisible lattices at
the terminal cut.
\end{proof}


\begin{proposition}[Finite-class colength in the long-orbit range]\label{prop:inert-long-colength}
Use the marked crystalline surjection and signed jump windows
of Lemma~\ref{hyp:long-orbit-data}, with the cyclic spacing
and arc exponents of Paragraph~\ref{par:inert-cut-exponents}.
Then \eqref{eq:long-orbit-exponent} gives a normalized colength on
integral finite classes at most \(3C_2U_B\), with \(U_B\) as in
Section~\ref{sec:cut-length-rows}.
There exists one rational place entry \(c_w\) for all good characters at \(w\), satisfying
\[
 3C_2U_B\le c_w\le(3C_2+1)U_B.
\]
\end{proposition}

\begin{proof}
The lifting argument behind \eqref{eq:long-orbit-exponent} is
in the two operators \(\Phi-1\) and \(\gamma-1\) of the Herr
complex. Here \(\Phi\) advances one letter and \(\gamma\) acts
within a letter. Work in the tame-invariant Wach coordinates,
with parameter \(X\) and Wach lattices \(N_i\) at the letters.
The first filtration estimate places \(\Phi(X^M N)\) in
\(X^M(\ell,X)N\) on the product \(N\) of the letters. On \(X^MN\) contraction therefore gives the convergent inverse
\[
 (\Phi-1)^{-1}=-\sum_{a\ge0}\Phi^a.
\]
The coefficients are complete and
\(\Phi\) preserves the ideal \((\ell,X)\).

\proofstep{Represent the scaled class by a restricted Herr pair and lift its arcs.}
A finite class in the target has a Herr representative attached to
a crystalline extension of the trivial lattice by \(T_2\). If \(E_0\) lifts the trivial generator in period coordinates, the target pair representing the connecting class is
\[
 ((\Phi-1)E_0,(\gamma-1)E_0).
\]
The
cut colength bound provides, at each cut \(c\), a vector \(x_c\)
in the extension Wach lattice over \(\ell^R\) times the trivial
generator. The difference \((\gamma-1)x_c\) belongs to \(XN_c\)
of the target.

In fact it belongs to \(X\) times the extension
Wach lattice by triviality modulo \(X\), and to the target
period module because its image in the trivial quotient vanishes.
Since \(X\) is regular on the trivial quotient and the
extension Wach lattice intersects the target period module in the
target Wach lattice, this gives the assertion. After advancing
\(j\) letters and multiplying by \(\ell^\sigma\), the correction
of Paragraph~\ref{par:wach-cut-correction} gives
\[
       Q_c=\ell^\sigma\Phi^j x_c-\tau_c,
       \qquad (\gamma-1)Q_c\in X^M N_{c+j}.
\]
Use the target Wach lattice as its image in the extension.

For consecutive cuts \(c',c\), write \(a\ge j\) for their
forward distance. At the letters from \(c'+j\) up to the
letter before \(c+j\), set \(E_{c'+j+t}=\Phi^tQ_{c'}\).
This gives a global vector \(E\) with trivial quotient
\(\ell^{R+\sigma}\) times the generator. The first component
\(\ell^\sigma(\Phi-1)E\) vanishes at all the letters within
each propagated arc. At a shifted cut it is
\[
 \begin{split}
   y_c&=\ell^\sigma(\Phi^a Q_{c'}-Q_c)
       =\ell^\sigma(\Phi^j w_c+\tau_c),\\
   w_c&=\Phi^{a-j}Q_{c'}-\ell^\sigma x_c.
 \end{split}
\]
The two terms defining \(w_c\) have the same trivial
quotient, so \(w_c\) is in the target. Since \(\gamma\)
commutes with \(\Phi\), the deep-lattice defects of the
\(Q_c\)'s propagate along each arc and
\((\gamma-1)y_c\in X^M N\).

Set \(z_c=(\Phi-1)^{-1}(\gamma-1)y_c\) there.
Uniqueness of the inverse on \(X^MN\) gives
\[
 \sum_c(y_c,z_c)
    =\bigl(\ell^\sigma(\Phi-1)E,\ell^\sigma(\gamma-1)E\bigr).
\]
The difference from \(\ell^{R+2\sigma}\) times the original
pair is the Herr coboundary of
\(\ell^\sigma E-\ell^{R+2\sigma}E_0\). This is a target
vector because its two terms have equal trivial quotient.

Let \(c'\) be the preceding cut and \(a\) the length of its arc
to \(c\). The propagated \(Q_{c'}\) at \(c\) belongs to the span
of \(\Phi^a\) of the extension Wach lattice at \(c'\), since its
correction at \(c'+j\) is in the span of \(\Phi^jN_{c'}\).
The signed bound into the Frobenius span in
Paragraph~\ref{par:wach-level-models} gives the same membership
after multiplying \(x_c\) by
\[
 \Lambda=\prod_{s<a}p_s^{H_s^{\mathrm{arc}}}.
\]

The filtration of this extension has the same bounds as the target
with the jump at zero included.

For \(\Lambda w_c\), membership in the extension Frobenius
span and in the target period module already holds. To
intersect these two modules, put \(\mathcal R=O[[X]]\)
and use a basis \((b_t)_t\) of \(\mathcal R\) over
\(\phi^a(\mathcal R)\). Such a basis comes from Weierstrass
division by \(\phi^a(X)-Y\) over \(O[[Y]]\), since
\(\phi^a(X)\) first has a unit coefficient in degree \(\ell^a\).
Group the extension expression as \(\sum_t b_t\Phi^a u_t\).
If \(d_t\) is the trivial quotient of \(u_t\), then
\[
       \sum_t b_t\Phi^a u_t\longmapsto
                    \left(\sum_t b_t\phi^a(d_t)\right)e_0=0 .
\]
Independence of the basis and injectivity of \(\phi^a\)
give \(d_t=0\) for every \(t\). Therefore every \(u_t\)
lies in the target Wach lattice at \(c'\). This proves
\(w_c\in W_2\). The two defects to be subtracted belong respectively to \(X^MN_c\) and \(XN_c\). Therefore
\[
 (\gamma-1)w_c\in X^M N_c+XN_c\subseteq XN_c.
\]


For a change of cut vector by \(u\in XN_c\), the correction
of Paragraph~\ref{par:wach-cut-correction} changes by
\(\ell^\sigma\Phi^j u\) modulo the transported error.
It cancels the change of the uncorrected advanced vector.
Hence the corrected first component changes by a vector
\[
 v\in X^M N.
\]
The cocycle equation determines the change
of the second component as \((\Phi-1)^{-1}(\gamma-1)v\).
Together these give the coboundary of \((\Phi-1)^{-1}v\).
Thus comparison of cut vectors modulo \(XN_c\) suffices
for the finite classes of the corrected arc pairs.

Choose a product \(D_{\rm prev}\) of cyclotomic factors
clearing \(\Lambda w_c\) into \(N_c\), by iterating the
one-step height inclusions on the preceding arc. Choose
\(D_{\rm adv}\) clearing \(\Phi^j N_c\) into \(N_{c+j}\)
and \(D_\tau\) clearing \(P_c\) into \(N_{c+j}\).
Then the product
\[
            D=\phi^j(D_{\rm prev}\Lambda)\,D_{\rm adv}D_\tau
\]
clears \(y_c=\ell^\sigma(\Phi^j w_c+\tau_c)\) into \(N_{c+j}\)
by semilinearity and \(\tau_c\in P_c\).
The second component \(z_c\) and the defect
\((\gamma-1)y_c=(\Phi-1)z_c\) are in \(X^MN\).
Thus the two components satisfy the actual finite
cyclotomic-denominator and deep-lattice tests.

Pass to the ordinary cyclotomic parameter \(\pi\). Put
\(V=T\otimes_O K\), with \(K=\operatorname{Frac}(O)\), and set
\[
 q_{\rm cyc}=\phi(\pi)/\pi,\qquad t=\log(1+\pi),\qquad
 \lambda=t/\pi .
\]
The factors \(p_s\) occurring in \(D\) are associates of
\(\phi^s(q_{\rm cyc})\). Also \((X^M)=(\pi^{(\ell-1)M})\).
Thus, for some finitely supported family of integers \(b_s\ge0\),
the Herr pair \((y,z)\) satisfies
\begin{equation}\label{eq:long-arc-herr-lattices}
 \Bigl(\prod_s\phi^s(q_{\rm cyc})^{b_s}\Bigr)y\in N(T),\qquad
 z\in\pi N(T),\qquad (\gamma-1)y=(\Phi-1)z .
\end{equation}
The integral Herr isomorphism gives a class in \(H^1(F,T)\).
Its rational image is represented by the same pair in the Robba-ring
Herr complex \cite[Theorems 2.3 and 2.8]{liu-herr}.
In the corresponding extension \(0\to V\to V'\to K\to0\), choose
a tame-invariant lift \(e_0\) of the trivial basis. The extension
actions are \(\Phi e_0=e_0+y\) and \(\gamma e_0=e_0+z\).

Let \(\mathcal B^+\) denote the analytic functions on \(|\pi|<1\).
Use crystalline coordinates in the kernel. For a positive crystalline
representation the Wach comparison and its elementary divisors put the
change of basis and its inverse over \(\mathcal B^+[1/\lambda]\)
\cite[Propositions II.2.1 and III.2.1]{berger-wach}.
A Tate twist gives the same conclusion for \(V\): the ratio of
the two twist generators is
\[
 \pi^{-k}/t^{-k}=\lambda^k.
\]

Since \(\lambda(0)=1\), the change of basis is regular at zero
in both directions. Writing brackets for crystalline coordinates gives
\[
 \lambda=\prod_{s\ge0}\frac{\phi^s(q_{\rm cyc})}{\ell},
 \qquad \lambda^B[y]\in(\mathcal B^+)^{\dim V}\quad(B\gg 0),
 \qquad [z](0)=0 .
\]
By \eqref{eq:long-arc-herr-lattices}, a sufficiently high power of the product absorbs all finitely
many cyclotomic denominator factors of \(y\) and the poles of the
change of basis. The factor \(\pi\) in \(z\) gives the last equality.

In the crystalline basis the equation \((\gamma-1)v=z\)
has a solution analytic near zero. Use \(t\) as local
coordinate and put \(\omega=\chi(\gamma)\). If
\(z(t)=\sum_{k\ge1} z_k t^k\), the solution with zero
constant is
\[
          v(t)=\sum_{k\ge1}\frac{z_k}{\omega^k-1}t^k .
\]
The valuations \(1+v_\ell(k)\) of the denominators grow
logarithmically in \(k\). Thus any strictly smaller positive disk of convergence of
\(z\) is a disk of convergence of \(v\). For \(u=y-\Phi v+v\)
the Herr equation and \(\gamma\Phi=\Phi\gamma\) give
\[
 (\gamma-1)u=(\Phi-1)z-(\Phi-1)(\gamma-1)v=0 .
\]

In the crystalline basis the coefficient \(u_k\) in degree \(k\)
is therefore killed by \(\chi(\gamma)^k-1\). For \(k>0\) this
factor is nonzero, so the coordinates of \(u\) are constant. For every \(m\ge1\)
the identity of germs is
\[
 v=\Phi^m v-\sum_{i=0}^{m-1}\Phi^i y+\sum_{i=0}^{m-1}\Phi^i u.
\]
On every disk of fixed radius less than one,
\(\phi^m(\pi)=(1+\pi)^{\ell^m}-1\) tends uniformly to zero.
Hence \(\Phi^m v\) is analytic there for sufficiently large
\(m\). Also \(\lambda^B\Phi^i y\) is analytic on the
whole open disk for every \(i\), since
\[
 \lambda/\phi^i\lambda=\prod_{k<i}\phi^k(q_{\rm cyc}/\ell).
\]

The displayed identity therefore extends \(\lambda^B v\)
to the whole disk. The continuations agree on overlaps by
their Taylor series at zero. It follows that \(v\) belongs
to the Robba module with \(t\) inverted. The identity
\[
 (\gamma-1)v=z
\]
holds there by its power series identity
at zero after clearing powers of \(\lambda\), so \(e_0-v\)
is \(\gamma\)-invariant. Its average over the tame
group is invariant under the full cyclotomic group and still
maps to the trivial basis. Berger's comparison
\cite[Theorem 0.2]{berger-differential} identifies these
invariants with \(D_{\rm cris}(V')\). The rank increases
by one from \(D_{\rm cris}(V)\), so \(V'\) is crystalline.
Thus the integral Herr class is finite. This criterion applies
to both lattices and is compatible with the same row map.

The quantitative estimate at an arc supplies a source vector
of this kind mapping to \(m_c w_c\) modulo \(XN_c\), where
\[
 m_c=u_c\ell^{2\sigma+2a_0+R+k_c},\qquad u_c\in O^\times,
                 \qquad 0\le k_c\le2\kappa.
\]
It applies the polynomials in \(\gamma\) to the target and then
to its lift to impose the required local exponent windows.
Modulo \(XN_c\) their actions reduce to their values at \(1\),
since \(w_c\) is \(\gamma\)-invariant modulo this submodule.

Let \(\omega\) be the cyclotomic value of \(\gamma\), with
\(v_\ell(\omega-1)=1\), and let \(E=a-1\le A\).
The polynomials have the form
\[
 P_{E,J}(Z)=\left(\sum_{v=0}^{\ell^E-1}Z^v\right)
            \prod_{s\in J\setminus\{0\}}
                 (Z^{\ell^E}-\omega^{s\ell^E}),
\]
At \(1\) the valuation is
\(E+\sum_{s\in J\setminus\{0\}}(1+E+v_\ell(s))\)
by the valuation of powers of a principal unit. The window
bound gives at most \(\kappa\) for each polynomial. Between the two
polynomials the comparison of Paragraph~\ref{par:wach-mixed-constant}
costs \(2\sigma+2a_0+R\) on the level inputs.
Paragraph~\ref{par:wach-cut-colength} gives the cut colength
on the two-sided Wach height bounds. The simultaneous level
realization and the two polynomial actions are those of
Paragraphs~\ref{par:wach-level-models}--\ref{par:wach-polynomial-actions}.

Multiply the finite source class for arc \(c\) by
\(u_c^{-1}\ell^{2\kappa-k_c}\), and sum over the cuts. Its image
is \(\ell^{2\sigma+2a_0+R+2\kappa}\) times the sum of the target
arc classes. Together with the initial decomposition this gives
\[
 e=2R+4\sigma+2a_0+2\kappa.
\]
The sum over the cuts costs no
additional power of \(\ell\).

\proofstep{Include the torsion generators and then take the normalized colength.}
The free part of the arc calculation bounds the failure of lifting along the orbit. Torsion in the generators must be included before the resulting length can enter the inert-prime table.


Let \(V_2=T_2\otimes_O\operatorname{Frac}(O)\) and
\(r_0=\dim V_2^{G_F}\). By the Bloch--Kato dimension formula the
free rank of the integral finite group is
\[
 \dim(D_{\mathrm{dR}}(V_2)/\mathrm{Fil}^0)+r_0\le dr+r_0.
\]

For a uniformizer \(\varpi\) of \(O\), the long exact sequence
for multiplication by \(\varpi\) gives
\[
 0\longrightarrow T_2^{G_F}/\varpi T_2^{G_F}
   \longrightarrow (T_2/\varpi T_2)^{G_F}
   \longrightarrow H^1(F,T_2)[\varpi]\longrightarrow0.
\]
The first term has dimension \(r_0\) and the middle term has
dimension at most \(r\). Thus at most \(r-r_0\) torsion generators
are needed. The rank of the free quotient and the number of torsion generators together cost at most
\[
 (dr+r_0)+(r-r_0)=(d+1)r.
\]
Since its quotient by the image of
the finite group of \(T_1\) is killed by \(\ell^e\), its length
divided by \(\operatorname{length}_O(O/\ell O)\) is at most
\((d+1)re\). This is the normalization used for the actual row
cokernel. It handles all coefficient ramification degrees.

Finally, for \(h\ge1\) use
\[
 \log_\ell h\le h/D_0+j+1,
\]
by splitting the logarithm at
\(D_0\) and using \(\log_\ell t\le t\) for \(t>0\).
When the word has a position with larger jumps, the bracket in the bound for
\(e\) is at most \(2H/D_0+2(G_P+1)(j+1)\). Otherwise \(h=g\le G_P\)
and it is at most \(3(G_P+1)(j+1)\).

In both cases the resulting
normalized colength is at most \(3C_2 U_B\), with \(U_B\) as in
Section~\ref{sec:cut-length-rows}. Choose a rational table entry \(u_B\) satisfying
\[
 3C_2U_B<u_B<(3C_2+1)U_B.
\]
The lower bound makes \(u_B\) valid for the cokernel at every good character. Such an entry exists since \(U_B>0\).
This gives case B of the ordinary congruence comparison on the
tied local data, with one constant and cutoff for every place,
coefficient field and marked character in the window.

Use the single word chosen for the large-orbit estimate.

Sum the arc bounds and convert the result to the rational entry in its marked row.
\label{cor:inert-long-orbit-1}

The filtration statement for the Wach module
of a crystalline representation over an unramified field is
Berger's filtered comparison \cite[Corollary III.4.5]{berger-wach}.
The exponent in \eqref{eq:long-orbit-exponent} belongs to the
submission's quantitative lifting argument.
\end{proof}

